\section*{Hoofdstuk \ref{ch:b2:mass-spec-atomic} --- Massaspectrometrie en atoomspectroscopie}

\begin{solution}{exo:b2:mass-spec-atomic:1}
78; $400/2 = 200$; $194 + 1 = 195$, $[\mathrm M + \mathrm H]^+$.
\end{solution}

\begin{solution}{exo:b2:mass-spec-atomic:2}
Oneven massa: een oneven aantal stikstofatomen, één voor een zo klein molecuul. \ce{C3H7NO} (propaanamide, 73)
of \ce{C4H11N} (butaan-1-amine, 73).
\end{solution}

\begin{solution}{exo:b2:mass-spec-atomic:3}
Bijna gelijk: één broomatoom. 3 : 1: één chlooratoom.
\end{solution}

\begin{solution}{exo:b2:mass-spec-atomic:4}
Geel: natrium, \qty{589.0}{nm} en \qty{589.6}{nm}. Karmijnrood: lithium, \qty{670.8}{nm}.
% ledger: asd:NaD2, asd:NaD1, asd:Li671
\end{solution}

\begin{solution}{exo:b2:mass-spec-atomic:5}
$n_C \approx 6.7/1.11 \approx 6$.
\end{solution}

\begin{solution}{exo:b2:mass-spec-atomic:6}
\ce{CO}: 27.9949; \ce{N2}: 28.0061; \ce{C2H4}: 28.0313 u. Oplossend vermogen: \ce{CO}/\ce{N2} $28/0.0112
\approx 2.5 \times 10^3$; \ce{N2}/\ce{C2H4} $28/0.0252 \approx 1.1 \times 10^3$; \ce{CO}/\ce{C2H4}
$28/0.0364 \approx 7.7 \times 10^2$.
% ledger: mass:C12, mass:O16, mass:N14, mass:H1
\end{solution}

\begin{solution}{exo:b2:mass-spec-atomic:7}
86: het \omterm{def:b2:mass-spec-atomic:molecular-ion}{molecuulion} van \ce{C5H10O}. 57: \ce{C2H5CO+}, een \omterm{def:b2:aromatic-substitution:friedel-crafts}{acyliumion} uit $\alpha$-splitsing (verlies van
een ethylradicaal, 29). 29: \ce{C2H5+}.
% ledger: ms:pentan-3-one
\end{solution}

\begin{solution}{exo:b2:mass-spec-atomic:8}
86: $\mathrm M^{+\bullet}$. 71: $\mathrm M - \ce{CH3}$, het \omterm{def:b2:aromatic-substitution:friedel-crafts}{acyliumion} \ce{C3H7CO+}. 43: \ce{CH3CO+}
(verlies van het propylradicaal). 58: \omterm{def:b2:mass-spec-atomic:fragmentation}{omlegging van McLafferty}, via een waterstofatoom op het
$\gamma$-koolstofatoom van de propylketen, met verlies van etheen. In pentaan-3-on hebben de ethylgroepen geen $\gamma$-koolstofatoom: geen
McLafferty-ion; zijn kleine piek bij 58 is de \ce{^{13}C}-isotopenpiek van het ion bij 57 (drie koolstofatomen,
ongeveer 3.3~\%).
% ledger: ms:pentan-2-one, ms:pentan-3-one
\end{solution}

\begin{solution}{exo:b2:mass-spec-atomic:9}
Een rechte door de oorsprong: helling $\sum c_iA_i/\sum c_i^2 = 1.553/30.0 \approx \qty{0.0518}{L/mg}$. Monster:
$0.128/0.0518 \approx \qty{2.47}{mg/L}$ koper.
\end{solution}

\begin{solution}{exo:b2:mass-spec-atomic:10}
$t = L\sqrt{m/(2eU)} = 1.00\sqrt{1000 \times 1.6605 \times 10^{-27}/(2 \times 1.6022 \times 10^{-19}
\times 2.00 \times 10^4)} \approx \qty{16.1}{\micro s}$. Voor 1001 u groeit $t$ met de factor
$\sqrt{1.001} \approx 1 + 0.0005$: $\Delta t \approx \qty{8.0}{ns}$.
% ledger: const:u, const:e
\end{solution}

\begin{solution}{exo:b2:mass-spec-atomic:11}
Twee chlooratomen: $(0.758 + 0.242x)^2$ geeft $0.575$, $0.367$, $0.0586$: $m/z$ 84, 86 en 88 in de verhouding
100 : 64 : 10.
\end{solution}

\begin{solution}{exo:b2:mass-spec-atomic:12}
\ce{^{40}Ar^{16}O+}: $39.9624 + 15.9949 = 55.9573$; \ce{^{56}Fe}: 55.9349; $R = 56/0.0224 \approx 2.5 \times
10^3$. Meet anders een ander isotoop van ijzer, \ce{^{57}Fe}, of verwijder \ce{ArO+} in een botsingscel
vóór de analysator.
% ledger: mass:Ar40, mass:Fe56, mass:O16
\end{solution}

\begin{solution}{pb:b2:mass-spec-atomic:1}
\textbf{1.} 156, 158 en 160 (met 157 en 159): het molecuul bestaat in verscheidene isotopische vormen.
\textbf{2.} Drie pieken met telkens twee eenheden ertussen, de middelste de grootste: één chlooratoom en één broomatoom.
\textbf{3.} 156.
\textbf{4.} Even massa: geen stikstof, of een even aantal stikstofatomen.
\textbf{5.} $156 - 79 - 35 = 42$: \ce{C3H6}.
\textbf{6.} \ce{C3H6BrCl}; onverzadigingsgraad $(2 \times 3 + 2 - 6 - 2)/2 = 0$: geen ring, geen dubbele
binding.
\textbf{7.} $0.5/14.5 \approx 3.4~\%$, $3.4/1.11 \approx 3$ koolstofatomen.
\textbf{8.} De piek bij 157 is gegeven tot op 0.1 op een schaal waarop hij 0.5 is: een onzekerheid van 20~\%.
\textbf{9.} \ce{C3H6^{81}Br^{35}Cl} en \ce{C3H6^{79}Br^{37}Cl}.
\textbf{10.} $3 \times 12 + 6 \times 1.007825 + 78.918338 + 34.968853 \approx 155.9341$.
\textbf{11.} $156/(156.151 - 155.934) \approx 7.2 \times 10^2$.
\textbf{12.} Ja: twee waterstofatomen minder, 154.
\textbf{13.} Er gaat een broomatoom (79 of 81) verloren: \ce{C3H6Cl+}, 77 met \ce{^{35}Cl}, 79 met \ce{^{37}Cl},
3 : 1.
\textbf{14.} Verlies van HBr: het radicaalkation $\mathrm{C_3H_5Cl^{+\bullet}}$, even massa.
\textbf{15.} \ce{C3H5+}, het allylkation, na het verlies van beide halogenen (Br, daarna HCl).
\textbf{16.} \ce{CH2Cl+} (49 en 51, 3 : 1).
\textbf{17.} \ce{CH2Br+} (93 en 95) en \ce{C2H4Br+} (107 en 109), elk 1 : 1: broom op één uiteinde van
de keten, het tweede ion door verlies van \ce{CH2Cl}.
\textbf{18.} De binding \ce{C-Br} is zwakker dan de binding \ce{C-Cl}, en \ce{Br} is het betere vertrekkende
radicaal.
\textbf{19.} Nee: geen pieken bij 127--131; \ce{CH2Cl+} en \ce{CH2Br+} tonen elk halogeen op een eindstandige
\ce{CH2}.
\textbf{20.} \ce{BrCH2CH2CH2Cl}, 1-broom-3-chloorpropaan.
\textbf{21.} 156/158/160: M; 77/79: M $-$ Br; 76: M $-$ HBr; 107/109: M $-$ \ce{CH2Cl};
93/95: \ce{CH2Br+}; 49/51: \ce{CH2Cl+}; 41: \ce{C3H5+}; 39, 27: \ce{C3H3+}, \ce{C2H3+}.
\textbf{22.} Nee: er is geen carbonylgroep.
\textbf{23.} 1-Broom-2-chloorpropaan, \ce{CH3CHClCH2Br}: verlies van \ce{CH2Br} geeft \ce{CH3CHCl+} bij 63/65.
\textbf{24.} \ce{C3H6^{81}Br^{37}Cl}.
\textbf{25.} $M$: $0.758 \times 0.5069 = 0.384$; $M+2$: $0.758 \times 0.4931 + 0.242 \times 0.5069 =
0.496$; $M+4$: $0.242 \times 0.4931 = 0.119$. Verhouding $\boldsymbol{\approx 3.2 : 4.2 : 1}$, ongeveer 3 : 4 :
1; het spectrum geeft $14.5 : 18.6 : 4.4 = 3.3 : 4.2 : 1$.
\end{solution}
